\begin{tikzpicture}[font=\small,>=stealth]
\draw[rounded corners=3pt,gray!45] (0,0) rectangle (14.2,12.8);
\node at (7.1,12.2) {同支时：两块倒数都为正，才能拼成固定总量 $3/8$};\node at (7.1,10.75) {例1：同支 $AF=3$，另一段 $BF=24$};
\draw[gray,dashed] (2.2,9.4) rectangle (8.2,10.299999999999999);\fill[blue!30] (2.2,9.6) rectangle (7.533333333333333,10.1);
\fill[orange!45] (7.533333333333333,9.6) rectangle (8.2,10.1);
\draw[violet,dashed] (8.2,9.35)--(8.2,10.45);\node[anchor=west] at (10.7,9.85) {正的余量};
\node at (5.2,9.049999999999999) {$1/3+1/24=3/8$};
\node at (7.1,7.199999999999999) {边界：$AF=8/3$，已占满总量};
\draw[gray,dashed] (2.2,5.85) rectangle (8.2,6.75);\fill[blue!30] (2.2,6.05) rectangle (8.2,6.55);
\draw[violet,dashed] (8.2,5.8)--(8.2,6.8999999999999995);\node[anchor=west] at (10.7,6.3) {余量为零};
\node at (5.2,5.5) {$1/BF=0$：没有有限另一段};
\node at (7.1,3.65) {例4：假设同支 $AF=2$，已超过总量};
\draw[gray,dashed] (2.2,2.3) rectangle (8.2,3.2);\fill[blue!30] (2.2,2.5) rectangle (10.2,3.0);
\draw[violet,dashed] (8.2,2.25)--(8.2,3.35);\node[anchor=west] at (10.7,2.75) {不可能有正余量};
\node at (5.2,1.95) {$1/2>3/8$：所设同支情形无解};
\node at (7.1,.55) {虚框宽度表示 $K=3/8$；色条画的是倒数，不是原线段长度。};
\end{tikzpicture}
